%
%      Brno University of Technology
%      Faculty of Information Technology
%
%      BSc Thesis	2006/2007
%
%      Railway Interlocking Simulator
%
%      Bedrich Hovorka
%-----------------------------------------------------
%
\documentclass[english]{cls/fitthesis}
\usepackage[latin2]{inputenc}
\usepackage{amsfonts}
\usepackage{amsthm}
\usepackage{url}

% packages amsmath, babel, graphics and epsfig (resp. graphicx) are already included

% =======================================================================
% package "hyperref" creates clickable links in pdf, if we use pdflatex
% problem is that package hyperref must be included last, so it cannot
% be in the template
\ifx\pdfoutput\undefined % not running under pdflatex
\else
  \usepackage{color}
  \usepackage[colorlinks,hyperindex,plainpages=false,pdftex]{hyperref}
  \definecolor{links}{rgb}{0.4,0.5,0}
  \definecolor{anchors}{rgb}{1,0,0}
  \def\AnchorColor{anchors}
  \def\LinkColor{links}
  \def\pdfBorderAttrs{/Border [0 0 0] } % without borders around links
  \pdfcompresslevel=9
\fi


%  routine section  --------------------------------------------------------
\newsymbol\blacktriangleright 1349

% for czechiso...
%\DeclareUrlCommand\url{\def\UrlLeft{<}\def\UrlRight{>} \urlstyle{tt}}

% insert image
\newcommand{\insertimage}[3][\textwidth]{
  \begin{figure}[htbp]
  \begin{center}
    \includegraphics[width=#1,keepaspectratio]{img/#2}
    \caption{#3}
    \label{#2}
  \end{center}
  \end{figure}
}

\newcommand{\insertplot}[2]{
  \begin{figure}[hbtp]
    \begin{center}
      \include{#1}
      \caption{#2}
      \label{#1}
    \end{center}
  \end{figure}
}

% :o) use only in math mode , TODO check?
\newcommand{\dv}{\Delta v}

\newcommand{\cfr}[1]{\texttt{#1}}

% reference (formatting) to source code
% TODO java-packages ?
\newcommand{\sourceRef}[1]{\cfr{#1}}

% next step - next state - math
\newcommand{\ns}[1][]{
  y^{#1}_{n + 1}
}

\newcommand{\nadepsanaStranka}[1]{
  \newpage
  \vspace* {1.75cm}
  \begin{flushleft}
    \begin{Huge}
{\bf #1}
    \end{Huge}
  \end{flushleft}
  \vspace* {2cm}
}

\newcommand{\zkratky}{
  \nadepsanaStranka{List of Used Abbreviations}
  \input{zkratky}
}

\newcommand{\beeeegin}{
  % generate title page
  \maketitle
  \pagenumbering{Roman}
  % after title page manually insert official assignment and license agreement
  \nadepsanaStranka{Assignment}
  \begin{enumerate}
    \item Study the functions and components of electronic railway interlocking used for traffic control in railway stations.
    \item Study simulation methods applicable for implementing a railway network simulator in a railway station, particularly discrete, continuous and combined simulation. Propose a suitable simulation approach, or combination of approaches for implementing this simulator. Propose a suitable data format for representing the railway network model.
    \item Design and implement a selected part of the simulator and demonstrate functionality on a sufficiently illustrative model.
    \item Discuss the obtained results and other possible directions of development.
  \end{enumerate}
  \medskip
  \begin{description}
    \item[Category:] Modeling and Simulation
    \item[Implementation language:] Java, or at your own discretion
    \item[Operating system:] Windows, Linux
    \item[Literature:] According to supervisor's instructions
    \item[Date of assignment:] November 1, 2006
    \item[Date of submission:] May 15, 2007
  \end{description}

  \nadepsanaStranka{License Agreement}
  The license agreement is stored in the archive of the Faculty
  of Information Technology, Brno University of Technology.

  \newpage
  \makeabstract\makedeclaration
  \clearpage
  \setcounter{tocdepth}{2}
  \tableofcontents
  \pagenumbering{arabic}
  \setcounter{page}{1}
}

%  global settings section ------------------------------------------

%Information about the work/project
%---------------------------------------------------------------------------
\projectinfo{
  %Work
  project=BP,           %type of work BP/DP
  year=2007,             %year
  date=\today,           %submission date
  %Title of work
  title.cs=Simulátor železničního stavědla,  %title of work in Czech
  title.en=Simulator of railway interlocking, %title of work in English
  %Author
  author=Bedřich Hovorka,   %name surname of author
  %author.title.p=Bc., %title before name (optional)
  %author.title.a=PhD, %title after name (optional)
  %Institute
  department=UITS,
  %Supervisor
  supervisor=David Martinek, %name surname of supervisor
  supervisor.title.p=Ing.,   %title before name (optional)
  %supervisor.title.a=~,    %title after name (optional)
  %
  %Keywords, abstracts, declaration and acknowledgment can be defined
  %either using the following parameters or using reserved macros (see below)
  %===========================================================================
  %Keywords
  keywords.cs={železniční stavědlo, kombinovaná simulace, XML, Java, OOP}, %keywords in Czech language
  keywords.en={railway interlocking, combined simulation, XML, Java, OOP}, %keywords in English language
  %Abstract
  abstract.cs={
  		 Stavědlo je dispečerské zařízení pro řízení dopravy. Dispečer určuje nastavováním
  		 výhybek a semaforů cestu vlakům.
  		 V této práci se zabývám návrhem a výstavbou základu simulátoru takového zařízení v jazyce Java.
  		 Nastudoval jsem funkce tohoto zařízení. Zabýval jsem se strukturou celé aplikace a
                 implementoval chování základních prvků.
	      }, % abstract in Czech language
  abstract.en={
                Railway interlocking is a dispatching facility for traffic control.
                The dispatcher controls train paths by setting switches and signals.
                In this thesis, I describe the design and implementation of a simple simulator of the facility
in the Java language. I have studied functions of this facility.
                I interested with structure of whole application.
                And I implemented behaviour of foundamental elements.
	      }, % abstract in English language
  %Declaration
  declaration={
    I declare that I have worked out this bachelor's thesis independently under
     the supervision of Ing. David Martinek.  %Additional information was provided by ........
    I have listed all literary sources and publications from which I have drawn.
    },
  %Acknowledgment (optional) TODO
  %acknowledgment={I would like to thank my supervisor Ing. David Martinek} % optional
}

% =======================================================================
% =================================================================
\begin{document}
\beeeegin

% ----------------------------------------------------------------

\chapter{Introduction}
  In this bachelor's thesis I dealt with the topic of traffic simulation.
  My work consisted of studying the system components,
  experimenting with existing implementations of railway traffic control games,
  finding a suitable simulation library and
  as the main activity -- designing and implementing the foundation of a simulator.
  First it was necessary to design the application structure, internal model representation and data format.
  Then I could implement a graphical interface for editing with the internal representation.
  Furthermore I mainly dealt with the simulation model.

  In the second chapter I present the modeled system.
  Its behavior is better shown on already existing programs.

  In the third chapter I mention numerical solution of differential equations. In the next section
  you will encounter basic concepts from modeling and simulation. I will briefly mention
  what discrete, continuous and combined simulation is, with combined simulation in somewhat more detail.
  %However, informal way of expression prevails here.

  In the fourth chapter I describe rather than commonly known properties of the used tools,
  %which can be easily obtained from the web, %sneaky
  my own observations and examples from the implementation. This applies mainly to sections about Java and XML.
  In the section about the \emph{jDisco} simulation library I will briefly explain how it is used, to avoid
  some misunderstandings, what I had to solve further in the implementation and what not.

  In the fifth chapter I selected, in my opinion, important ideas from the design and implementation description.
  I go into details in case I find it necessary.
  %Greater emphasis on detail could degenerate into a dry and pointless enumeration.
  First I reflect on the overall application structure. Then I describe the organization of elements
  in the system. Followed by a description of how simulation is implemented. And finally
  I present an example demonstrating what was implemented.

  In the final chapter I go through some possible extensions, especially those that closely follow the current implementation.
  This section is followed only by an evaluation of the entire work.

Some terminology was used from \cite{z173_1995sb} and existing implementations.
Also names in figures and schemes will be mostly in English to match the source code.
For understanding this text I assume that the reader is at least an experienced computer science student
and understands concepts from discrete and numerical mathematics, algorithms and object orientation.
%The text will be written concisely with cross-references between chapters on connections without unnecessary verbosity.

Example of reference to source files: \sourceRef{Example}.
\section{Stages and Goals}

    \subsection*{What was implemented before the bachelor's project}
    During the Semester project I created a prototype of an object model of the network, especially static properties.
    Based on this I designed an XML format and implemented a factory for loading and saving.
    Within the Modeling and Simulation project I created a combined model of train movement
    according to a given sequence of semaphores and edges (not in a network).
    For the User Interface Design course I submitted a simple graphical editor of the aforementioned network model,
    which displays the network in a familiar scheme similar to that used in electronic interlockings.
    I completed everything during the bachelor's project to the following goals:
    \subsection*{Goals of the bachelor's project}
    Create the core of a portable program:
    \begin{itemize}
      \item based on studied materials, design an extensible system of classes, connectable communicating modules
      \item Dynamic properties and behavior of basic elements -- semaphore, switch, track
      \item Movement and localization of trains in the network
      \item Data structure for manipulating paths, especially setting paths (without their cancellation)
      \item Create an example that demonstrates the functionality of the aforementioned goals
    \end{itemize}
    I would leave modeling of other elements and all system details as extensions.

\chapter{Modeled System and Existing Implementations}
This chapter will briefly describe the functions of the modeled system. I will also list some already existing applications.
Facts that need to be taken into account will also be mentioned here, even if their implementation was not completed.
  \section{Railway Interlocking}
  Railway interlocking is a device that allows controlling a traffic network of some node
  or track section from one place. It controls the behavior of several elements.
  For example, it sets signals on semaphores and switches rail switches.
  Some today's interlockings have a computer interface and are operated by one dispatcher. He allocates
  a certain path to each train by entering the initial and target semaphore. It is then up to the interlocking to find the most suitable
  unoccupied path.
  \begin{enumerate}
    \item If it doesn't find it due to all possibilities being occupied, it saves the request (pair -- start and destination) to a queue.
          When releasing tracks it tests the possibility of assembling the first element in the queue. The queue can be edited.
    \item Otherwise it sets elements in the path
  \end{enumerate}
  This is also related to safety. The system must not allow more than one
  train path on each track\footnote{There are also so-called shunting paths, but I don't consider them yet}.
  The following section will reveal what such a simulator can look like.

  \section{Existing Implementations}\label{recenzes}
  For a better idea of how such an application works, I describe here
  some already existing railway traffic control simulators. Although they are mostly games, in their concept
  they are close to a real system trainer. However, they maintain simplicity of operation.
  I also used them as a means for system analysis. Experimenting with these programs as a method
  of reverse engineering\footnote{\url{http://en.wikipedia.org/wiki/Reverse_engineering}} allows studying the structure
  and behavior of most of the system. Each has its national specifics (details in control rules), but many principles are common, which I describe for the first one. For the others I only mention differences.
    \subsection*{Staničář}
      This game simulator is intended to model Czech interlocking systems.
      The user chooses time and station. Stations and trains are stored in one large XML file.
      Then the track layout and list of trains in the system are displayed.
      Now they can control individual elements and observe system behavior.
      Trains enter the track layout at designated places (according to timetable, generator with exponential
      distribution of time between generations).
      The task is to guide trains to their destination according to timetable and maintain traffic flow.
      Each station or larger system has its stack (here called FIFO queue) of commands, where
      they are stored when they cannot currently be executed. For each train there is an option to view the timetable
      (time, stop, track number). Stations can be created in a separate editor.
      There is also a timetable editor.
      In the current version this game runs under .NET 2.0 in Windows. Earlier versions were executable in dotGNU.
    \subsection*{SimSig}
      This British simulator supports multiple stations, but each has its own installer.
      Functionally it is more comprehensive. It can simulate power outages, random disruptions in various parts of the network,
      and other failures of trains and network elements. There is also
      a certain degree of automatic control even in branched track layouts (timetable execution).
      Although intended for Windows, unlike the previous one I managed to run it
      in wine\footnote{wine-0.9.36}.
    \begin{description}
      %\item[TrainDispatcher] United States
      \item[Gordikon, Multikon, Brno] The network is fixed.
        Unlike the others, at the beginning of simulation trains are also placed on tracks in the network,
        not just that they must first enter the system through a special element in the network, which is more realistic.
    \end{description}
\chapter{Theory}
In this chapter I describe some theoretical foundations. Terms like \emph{set} or \emph{relation} I assume the reader knows.
    \section{Numerical Methods}\label{numMethods}
      For computing differential equations on digital computers, numerical methods are usually most suitable,
      because operations are performed discretely here. Continuous computations must be converted to them.
      And numerical methods are precisely this conversion means.
      Depending on usage, it is necessary to examine their properties such as accuracy and stability of solution.
      This problem is dealt with in more detail in \cite{rektorys}. By decreasing the step size, computation accuracy increases
      (differently for different methods), but at very small steps, the imprecision of elementary operations of digital arithmetic plays a role.

      % basic principle NRODR, step change,
      As an example I present here a single-step method for solving ODE with initial conditions -- \emph{Runge-Kutta-Fehlberg}
      (equations \ref{eqRKF45_beg} - \ref{eqRKF45_end}), because it is used further in this thesis. It is of fifth order and
      moreover we can combine the intermediate values $k_n$ into a fourth-order approximation (equation \ref{eqRKF45_4aprox}). From that we can
      then calculate an error estimate $ \ns - \ns[*] $ . Use of these properties see section \ref{modelling}. In figure
      \ref{rkf45} we see its use when we expect a result close to a second-degree polynomial.

      \begin{eqnarray}
        k_1 & = & f(x_n, y_n) , \label{eqRKF45_beg} \\
        k_2 & = & f(x_n + \frac{1}{4}h, y_n + \frac{1}{4}hk_1) , \\
        k_3 & = & f(x_n + \frac{3}{8}h, y_n + \frac{1}{32}h(3k_1 + 9k_2)) , \\
        k_4 & = & f(x_n + \frac{12}{13}h, y_n + \frac{1}{2\,197}h(1\,932k_1 - 7\,200k_2 + 7\,296k_3)) , \\
        k_5 & = & f(x_n + h, y_n + h(\frac{439}{216}k_1 - 8k_2 + \frac{3\,680}{513}k_3 - \frac{845}{4\,104}k_4)) , \\
        k_6 & = & f(x_n + \frac{1}{2}h, y_n + h(-\frac{8}{27}k_1 + 2k_2 - \frac{3\,544}{2\,565}k_3 +
                          \frac{1\,859}{4\,104}k_4 - \frac{11}{40}k_5)) \\
        \ns & = & y_n + h(\frac{16}{135}k_1 + \frac{6\,656}{12\,852}k_3 + \frac{28\,561}{56\,430}k_4
                        - \frac{9}{50}k_5 + \frac{2}{55}k_6) \label{eqRKF45_end} \\
        \ns[*] & = & y_n + h(\frac{25}{216}k_1 + \frac{1\,408}{2\,565}k_3 + \frac{2\,197}{4\,104}k_4 -
                          \frac{1}{5}k_5) \label{eqRKF45_4aprox}
      \end{eqnarray}

      \insertplot{rkf45}{Demonstration of method step}

  %\section{Object-oriented modeling}
  %... \cite{umlDesign}
  \section{Modeling and Simulation}\label{modelling}
    In this section I present basic concepts from modeling and simulation.
    Details in \cite{imsSkripta, helsgaun2001}
    \begin{description}
      \item[System] is anything for which we can describe its behavior. It can be real or imaginary. For example
 	a railway network with trains and control elements.
      \item[System element] is an elementary, further indivisible, part of the system. Mutual interaction of elements determines behavior
        of the entire system. For example a track section. Here it depends on the level of abstraction.
      \item[Model] is a system that mimics properties and behavior of the original system.
      \item[Modeling] is the gradual creation of models. From observation and study of the system arises
      a \emph{conceptual model}, which is a set of uncertain information (outline).
      From it we create an \emph{abstract model} -- formal description of the system.
      And from that by isomorphic mapping we program a \emph{simulation model}.
      \item[Simulation] Acquiring new knowledge about the system by experimenting with the simulation model.
      \item[Real time] actual time in which the model runs
      \item[Model time] time axis of the model
      \item[Machine time] time that was needed for computation (i.e. how long the processor was switched to the context
       of the given process)
    \end{description}

    \subsection{Discrete Simulation}
      The state of all elements is defined only at discrete time moments.
      During this moment an atomic operation called an event is performed, which can change this state.
      For formal description, Petri nets are often used. Simulation is controlled using an event calendar
      with this simple algorithm:

      \begin{verbatim}
      	for event in calendar.remove_iterator :
      	  model_time = event.time
      	  event.behaviour()
      \end{verbatim}

      %\subsubsection{Petriho síť}
      %Petriho síť je pětice $(P,T,F,M_0,W)$ kde:
      %\begin{itemize}
      %	\item $P$ je množina míst
      %	\item $T$ je množina přechodů
      %	\item $F \subseteq (S \times T) \cup (T \times S)$
      %	\item $
      %\end{itemize}

    \subsection{Continuous Simulation}
      The state of each element is defined at every moment of the specified time interval.
      The abstract model consists of algebraic, difference or differential equations and their systems.
      For single-step methods, based on error estimate the step can be changed and thus accelerate computation.
      The method presented in section \ref{numMethods} has such properties.

      \begin{verbatim}
        #invalidate previous changes (derivative=0)
        for variable in all_variables : variable.state_fix()
        #compute state changes from previous
        for continous in all_cont_processes : continous.derivatives()
        integrate() #step of numerical method with model time shift
      \end{verbatim}

    \subsection{Combined Simulation}\label{sim:comb}
      The system contains both continuous and discrete elements. Simulation control must during computation of continuous states
      change the step to reach the planned discrete event. Furthermore it must detect state conditions and
      schedule execution of the event assigned to the condition.
      From numerical methods, single-step ones are therefore suitable, because for them only the previous result matters.

      \subsubsection{Petri net when creating a combined model}
      During application design I used Petri nets for clear
      expression of processes. This way I better imagined the simulation flow and how to express it in programming
      language. Mostly these were C/E nets, i.e. a place, here called condition, contains at most
      one token and a transition, here called event, is executable in case all conditions before it hold and conditions after it do not hold. Depending on better expression I
      declared transitions either as events reacting to a state condition, which cancels
      validity of the token before it, or as a continuous process and the place after it is a state condition, which the
      following event waits for.

      Among other things I could also verify them using the CESim tool (see \cite{novosad2006}),
      if the net could be expressed without sacrificing clarity.
      For example the net in figure \ref{autoblok} is verified by this tool.

      \insertimage{trainProcess}{Petri net of train control according to signal aspects}

      In the net in figure \ref{trainProcess} there are additionally external conditions and priority.
      The process returns to the same state -- i.e. is before a signal
      , when the train arrives at the next signal or by signal change
      waiting before it ended. Conditions \uv{red} and \uv{allowing signal} represent the state of this signal
      and of course always\footnote{at each discrete moment of process state} exactly one of them holds.
      With condition \uv{$v \le 0$, next red} and priority I handle the possibility of stopping at the end of the following block,
      which was preceded by acceleration at its beginning.

\chapter{Used Languages and Tools}
  This chapter describes languages or tools that I used in my work.
  I became familiar with Java and XML earlier, so I don't want to go
  into details. On the other hand I will write here my opinions,
  why I used these languages or tools and possibly for what they were unsuitable in my work.

  \section{XML} \label{secXML}
  XML is a standardized markup language for storing various types of data. It extends text with marks, called tags, which
  are enclosed in $<$ $>$, have a name and can have several attributes. Tags are opening (starting with $<$), closing
  (starting with $</$) or ``self-closing'' (opening with $/$ before $>$).
  It is used in various areas, such as web pages, vector graphics, real estate registry, object
  serialization (in JavaBeans) and of course also for storing models in M\&S (e.g. Ptolemy II.).

  XML defines a general syntax in which it is possible to create custom tags.
  The internal representation is an n-ary tree of various node types like element, attribute, text etc.
  Data can thus be arranged hierarchically according to logical relationships.
  It is easily extensible and is supported by many languages, which allows creating data formats gradually by prototyping.
  Thanks to namespaces it is possible to
  insert into a document of one type a part of another (perhaps already existing). For example into a traffic network description, which
  represents mainly topology, one can add to each path its geometric properties using an external vector
  graphics format. For machine data exchange it is an excellent language, but for manual writing or even programming
  (e.g. XSLT) it seems unsuitable to me. Often it is advisable to compress it before sending over network, because here manifests
  the disadvantage of this type of markup -- repeated information overhead.

  \subsection{XML Parsers}
  \paragraph{Document Validation} Lexical and syntactic rules are the same for all XML.
  The creator of an XML application determines the semantics. The following definition tools can partially help.
  For determining how the so-called \emph{well-formed document} should look
  in XML, Document Type Definition (DTD) serves,
  where metadata is determined -- what tags the document can contain, what attributes they have and what they can
  contain, but only at the level of XML elements, missing for example type checking of attributes etc.
  Thus DTD was not enough and XML Schema \cite{XMLSchema2004} had to be created,
  which is an XML application describing additionally also object relationships between tags.

  In favor of choosing XML as data format speaks also that
  parsers already exist for various languages
  %For example in Java they are even part of the standard library (from version 1.4)}.
  They take care of lexical and syntactic
  analysis according to the general XML grammar and possibly perform validation.
  There are two main ones: SAX and DOM. SAX walks through the document and calls methods of a registered \cfr{Handler}, when
  it encounters the beginning or end of an element. During this walk validation can also be performed. In DOM parsing,
  SAX is used to convert text to the standardized internal form of the document (tree).
  Therefore it is used rather in programs
  that need random access to various levels of the document, such as editors and browsers. However, if an application has
  its own better internal data representation or long sequential transformations need to be performed, SAX is more suitable. More about XML in \cite{nutshellXML}

  \section{Java} \label{java}
  Java is a general object-oriented programming language.
  Its main advantage is that it runs in a virtual machine environment, whose implementation
  is available for many common operating systems. Thus it allows creating portable
  executable files.
  I chose it because I have the most experience with this language.
  For implementation I finally chose version 6 Standard Edition.
  There is also a compiler in GCC but it is incompatible with the latest official version.
  Moreover from version 7 the official version should also be open (see OpenJava project).
  Many publications exist -- e.g. \cite{pecinovskyJ5u, spellProfJava}.
  Java since its inception has undergone many innovations, much has been added and some library classes
  are deprecated. Therefore some information in older literature is outdated.
  The most reliable source is thus \cite{javaapidoc}.
  In Java I am used to writing code so that variable and method names are self-commenting.

    \paragraph{Containers (Collections)} Are analogous to STL from C++ and serve for storing a predefined amount of data. In Java these are objects, thus they themselves can be elements of another collection.
    For example \sourceRef{Map} serves for creating mappings from object to object, which reduces
    creating references directly in objects and thus also dependency of code of one class on another.
    The issue of containers is elaborated more in \cite{heroutUtil}, but when using Java 5 it is
    necessary to take into account that containers are parameterized.
    By composing collections I created custom data structures. Primarily I supplemented or fully
    defined interfaces and implemented unoptimized prototypes. From them can evolve storage
    tailored directly, yet they can still maintain a certain degree of generality. For example \sourceRef{Array2DMap}
    in speed is comparable to \sourceRef{TreeMap}, but
    most frequently called is just item retrieval and is performed during critical computations.
    Here it was mainly about minimizing creation of key objects of integer pairs while maintaining
    compatibility with \sourceRef{Map}.

    \paragraph{Advantages and Disadvantages}\label{java:col}
    By its syntax it inclines more to general languages (C, C++), yet by amount of standard libraries it catches up with many scripting languages.
    Strict static type checking may at first glance delay development,
    but with proper design and use of suitable editors it conversely speeds up debugging, because I can
    concentrate mainly on logical errors in code.
    What however unlike scripting languages it lacks is support for functional programming.
    For modeling complex systems with structures with recursive aggregation
    (e.g. path, which is a subclass of track, split into parts, which are also subclasses of track)
    for determining a message sent to all parts from the path I would rather implement using delegation and method references.
    Here namely a disproportionate amount of similar (copied) code can arise, which I partially
    solved using reflection, but this solution is inefficient, because the method must be found by string with name and those strings must be \uv{registered} additionally.
  \section{JDisco}
  Implementing simulation routines and additionally model code is a longer-term matter.
  Many simulation libraries exist\footnote{SimPack, JDEVS, then several identically named JSims and JavaSims},
  but mostly they are only discrete. These libraries I could extend with combined simulation.
  But why, when there already exists a ready, quite easily usable
  library. Moreover I fear that this implementation by itself would be worth at least one separate
  project. Another possibility is connection through native methods
  with a library in another language (e.g. SIMLIB). That is also time-consuming.

  \emph{JDisco} -- \cite{helsgaun2001} is a package of Java classes for description and simulation of combined models.
  It was developed at Roskilde University in Denmark.
  There are two types of processes: \emph{discrete} and \emph{continuous} (as was said in \ref{modelling}).
  Between scheduled \emph{discrete} processes are inactive phases, in which \emph{continuous} ones can be active.
  For description of continuous variables serves the class \sourceRef{jDisco.Variable}. It has two main attributes:
  instantaneous value -- \cfr{state} and instantaneous derivative -- \cfr{rate}.
  Continuous processes I create by inheriting from \sourceRef{jDisco.Continuous} and mandatory implementation of method
  \cfr{derivatives}, in which relationships between variables are described.
  Both these classes have methods \cfr{start} and \cfr{stop},
  by which I can determine when continuous computation should run by calling them in event description.
  In other words it is possible to define intervals of continuity of processes or variables this way.
  By inheriting from \sourceRef{jDisco.Process} and implementing method \cfr{actions} a discrete process arises.
  By implementing interface \sourceRef{jDisco.Condition} I create a state condition, which the process will
  wait for in calling method \cfr{waitUntil}. Furthermore the package contains other classes for simplifying model description,
  also several numerical methods (various \cfr{Monitors}). I will not make their list here.
  Documentation of this library I added\footnote{resp. it generates with program documentation}
  to program documentation of my thesis. More thus there.

    \subsection{Simulation Control}
    Simulation is controlled in background by chosen \cfr{Monitor}, which is responsible for:
    \begin{enumerate}
      \item model state changes continuously between events
      \item discrete events execute at correct time
      \item collection of information about model behavior is performed
    \end{enumerate}
    Continuous processes work in parallel and are synchronized with discrete ones, which are performed using calendar
    quasi-parallel. At one moment of simulation thus runs either one discrete process or \cfr{Monitor}
    controls computation of several continuous ones.
    Process is a Java thread giving up the processor just before
    calling \sourceRef{Object.wait} on itself, calls \sourceRef{Object.notify} to process following in calendar.

\chapter{Design and Implementation}
This chapter focuses on a brief description of implemented components. Details are contained in program documentation and source files. First the overall structure of source codes is described and then individual components are elaborated, primarily the most important ones.

  \section{Division into Modules -- Packages}
  First of all it is necessary to realize and keep in mind the overall application structure, even though only
  some parts will be implemented from it for now. It could happen that I implement something and then will laboriously
  redo it\footnote{during design it is not possible to keep in mind all details, especially where requirements change --
  so here and there a smaller rework is done}, because it won't connect.
  In figure \ref{dataFlow} are shown at higher abstraction basic functional modules of the application and communication between
  them, as I assumed at the beginning of design.
  \insertimage[300pt]{dataFlow}{Basic application structure and data flow}
  The program should work in two modes: editing and simulation. Here I decided that it would be better to create a wrapper
  for access to the data model. Modes have common operations, but also those that are possible only in that
  certain mode. This thus encapsulates the interface\footnote{Java interfaces -- in section \ref{java}}
  \sourceRef{Context} and its sub-interfaces.
  Its only implementation so far \sourceRef{DefaultContext} at first glance looks like an all-knowing class
  or like a storage for unplaceable functionalities. But the fact is
  that here are mostly concentrated methods connecting various parts of the program.
  For project foundation however I don't need to have two different implementations yet, it suffices to mind separation using interfaces. But it allows that simulation context won't be only in memory, but can be shared between applications
  for example over network or in database.

  \insertimage[211pt]{packageDependencies}{Graph of direct package dependencies}
  Between modules (packages) dependencies arise, i.e. code of one package refers to another.
  Especially unpleasant are cyclic ones, which are not a sign of
  good design, if they arise uncontrolled in large amounts. If there is an effort to create modules as independent as possible,
  code will be more reusable. In figure \ref{packageDependencies} are shown
  established module dependencies in current implementation.
  I assume that transitive ones will follow for the reader from it.
  This principle I held except for some exceptions, for which I will probably pay in future.
  Backward dependency of data on context is however solved through interface. Moreover for example class \sourceRef{AbstractPath}
  is rather a foundation for control package.

  All packages from figure \ref{packageDependencies} are also dependent on package \sourceRef{util}, which
  is dependent only on standard libraries. Especially on collections, which it extends with other needed data
  structures, such as undirected graph. About these structures more on page \pageref{java:col}.
  In following subsections I will describe other packages individually.

  \section{Network and Data Model}
    In this section is described in more detail the internal representation of the network.
    \cite{algorithms}
    Partially I was also inspired by \cite{Fleischanderl2001}.
    %  Object preserves element state, thus ensures it is defined at discrete time moments.

    From dispatcher's viewpoint it is an undirected graph $ G = (U, H) $.
    To each pair of objects $(node1, node2)$ where $node1 \neq node2$ is assigned at most one edge object.
    This graph (primarily its nodes)
    is placed in two-dimensional space with integer indices. This space is a kind of pseudo-raster
    -- matrix of objects. Edge is divided into parts and they fill free cells
    on a line\footnote{Bresenham's algorithm} between them. In each cell can of course be located at most
    one object.
    Nodes are connected to edges by determining cell neighborhood (found during editing), which encapsulates
    enumeration type \sourceRef{Cell.Segment}. Each cell has a set of segments\footnote{also serves for rendering}
    , and only to this neighborhood can be connected to each pair $(segment, node)$ at most one edge.
    From train's viewpoint the block, represented by edge, can consist of multiple sections, between which are by dispatcher
    uncontrollable (automatic) signals. Furthermore there exists so-called \sourceRef{TrackFacility}
    -- it is the smallest even multi-section unit, in which can be located at most one train.
    \insertimage[300pt]{trackLifeCycle}{Life cycle of \sourceRef{TrackFacility}}
    In figure \ref{trackLifeCycle} are shown track states and what can trigger their transition.
    These states are discrete and I expressed them using enumeration type. It should be clear that a train
    can enter only a reserved track, also that reserving is not possible if it is already reserved or occupied.
    It is reserved by dispatcher command and occupied resp. released by train.
    \subsection{Network Elements}
      \begin{itemize}
        \item Nodes: \sourceRef{NodeCell} -- I assume that in each node is a sensor,
          influencing occupation of connected edges.
          Its activity is simulated by train movement (see \ref{trainMoveSection})
          Except for switch, following nodes are directed.
          \begin{description}
            \item[Switch] is set either to main direction or to diverging
            \item[Main signal] changing signal aspect
            \item[Entry-exit point] represents external track system.
                 Or they can be as extension connected two inside model without simulation Worker
                 -- support for grade-separated crossings.
            \item[Track end] constant signal \uv{Stop} %until however turn of train -> dead end
          \end{description}
        \item Edges: \sourceRef{TrackBlock}
          Figure \ref{trackDiagram} illustrates derivation of block from track.
          One of possible edges represents simple track \sourceRef{SimpleTrackBlock}.
          It is \sourceRef{TrackFacility} with one section.
          Other edges can be inter-station sections:
          \begin{description}
            \item[Automatic block] has multiple sections -- see page \pageref{autoblokSection}
            \item[Semi-automatic block] facility, 1 main section + 2 pre-station
          \end{description}
          \insertimage{trackDiagram}{Diagram of track derivatives}
      \end{itemize}

    \subsection{Dynamic Configuration and Paths in Network}
      Reserved tracks and allowing signals represent one or more so-called
      "set train routes".
      This path gradually dissolves by train movement and can change under certain conditions.
      Formally it is a path from graph theory i.e. sequence of nodes and unoccupied edges,
      but with regard to element configuration -- more specifically:
      \begin{itemize}
        \item initial and final node is a signal in path direction or entry-exit point
        \item switches connect edges (both neighbors from sequence)
        \item signals in direction must allow movement, except last one, which forbids it
        \item opposite direction signals: main ones forbid (further in not-implemented -- section ones are off and on distant signal path has no effect)
      \end{itemize}
      This functionality\footnote{working with main signals, because others are inside block}
      is implemented in \sourceRef{AbstractPath}. In figure \ref{trackDiagram} is also
      shown connection of path to track.

    \subsection{Data File}
      For storing and loading data I designed so-called \sourceRef{ContextFactory}.
      I chose XML with validation using XML Schema (see \ref{secXML}).
      Current data format is in basic form and represents assembly of network elements at lowest
      user-editable level. It represents static description of station network.
      Individual nodes and edges are simply enumerated as sub-elements of root element.
      For distinguishing element types I use among others enumeration types. Not only do they save space and are easily extensible,
      but they are also easily loaded and written.
      %In section \ref{chapExtensions} I will deal with possible extensions.

  \section{Simulation}
    This section describes simulation flow -- by informal description and using selected
    abstract models of most important processes.

    \subsubsection{Why Combined Simulation?}
      Train position can be conceived discretely - i.e. which track circuits it occupies.
      However it is also necessary to know its more precise position, speed and acceleration at each moment.
      The degree of precision is adjustable. It allows more possibilities and changes better. For example
      one can change equation in motion description.

    \subsection{Entry-exit Points -- \sourceRef{InOut}}
      Represents entry and exit to track layout.
      Its behavior, which is based on train queue, describes \sourceRef{InOutWorker}.
      In simple models one can insert into it directly by generator,
      in more complex ones higher control logic must do it -- dispatcher.
      By itself it can allow path to nearest signal oriented in direction of movement, if
      switches connect some path\footnote{here it is possible to shield entry}
      and at the same time there is some train in queue and at the same time all tracks on it are free.
      In future it can represent connection of external track layout.
      Currently I consider models with two \sourceRef{InOut}.
      When I connect only two these \sourceRef{InOut} with one track, a competitive
      fight for that track breaks out between them. Whoever comes first can let train into system. However it cannot
      happen that both do it at once or let trains against each other.

    \subsection{Train Movement} \label{trainMoveSection}
      Train waits in entry-exit point queue for movement permission. After obtaining it
      subprocess \sourceRef{Train.Front} starts, which when encountering a node
      (\sourceRef{PathSeparator}), possibly changes parameters according to it and occupies section before it.
      When train exits completely \sourceRef{Train.Tail} starts. It when passing through each node
      releases section behind it.

      Each signal provides information about permitted speed behind it and about signal aspect at
      following signal, unless signal \uv{Stop} is showing. In this case the
      train should brake before it and wait for signal allowing movement.
      For simplification I assume that signal state can be detected only in discrete state,
      when train arrives at signal, because determining when it can register signal
      change is in reality indeterminate unlike in program.

      %If it is at station, at
      %which it has scheduled departure time, it must stop here and wait for scheduled time and at least
      %until passengers board and alight (waits for them with normal distribution with mean
      %2 and variance 1 minute), then it can depart.
      When passing by signal it starts changing speed depending on signal aspect.
      So-called "start on caution", when train at first signal starts accelerating and before
      following one stops, is handled separately. Subprocesses end at terminal
      entry-exit points \sourceRef{InOut}. Here I simplified model by assumption
      that maximum train length must surely be less than minimum possible length of track block between two
      \sourceRef{InOut}. Such similar edge cases are not worth handling otherwise than declaring model
      erroneous. Detection of this error requires finding minimal paths.

      Formally it is a combined process (see section \ref{sim:comb}). Discrete part waits for state conditions
      of arrival at signal and sets acceleration. This interaction describes Petri net in figure \ref{trainProcess},
      it serves however only for illustration. Simulation is not implemented directly using PN.
      I chose simple continuous model, because from viewpoint of debugging discrete interactions
      it is more predictable.
      Basic motion equations \ref{vaIntegrator} and \ref{svIntegrator}
      for acceleration $a$, velocity $v$ and distance $s$:
	\nopagebreak
        \begin{eqnarray}
        v & = & \int{a\, \mathrm{d}t} \label{vaIntegrator}\\
        s & = & \int{v\, \mathrm{d}t} \label{svIntegrator}
        \end{eqnarray}

      Computation of acceleration $a$ in equation \ref{aInSpeedChange}
      from initial velocity $v_0$, target velocity $v_1$ and distance $s$ for acceleration:
      \nopagebreak
      \begin{eqnarray}
         \dv & = & v_1 - v_0 \nonumber \\
         a  & = & \frac{\dv (v_1 + v_0)}{2s} \label{aInSpeedChange}
         %t  & = & \left| \frac{\dv}{a} \right| \label{tInSpeedChange}
      \end{eqnarray}

      \section{Example -- Shunting Loop}
      For better demonstration how simulation model works I created a \uv{shunting loop} controlled based on simple rules.
      To each controllable signal is permanently assigned at least one created path, which leads to another
      track either with possibility of stopover or leaving system,
      thus from certain viewpoint represents one indivisible transfer operation.
      Trains are created with simple timetable (two data from where, to where -- \sourceRef{InOut}) using generator
      with exponential distribution of time between individual generations.
      Before they are inserted into \sourceRef{InOut} queue, congestion of track layout should be prevented.
      Control process iteratively goes through knowledge about system:

      \begin{enumerate}
        \item forgets trains that already passed through system
        \item from trains waiting for entry selects enough so that at most two are in system in total
              and grants them permission (by activating process)
        \item goes through all track blocks
	      \begin{enumerate}
	        \item in which is located a train or reserved path, determines its resp. its direction
                      -- signal to which train is guided
	        \item if from this signal is free path (previously stored) sets it
              \end{enumerate}
      \end{enumerate}
      \insertimage{vyhybna}{Track diagram of shunting loop}
      To diagram in figure \ref{vyhybna}: line segments are tracks, their connections are switches,
                   $\bullet$ is \sourceRef{InOut} and~$\blacktriangleright$ is signal.
      After simulation start messages are printed line by line to standard program output:
      \begin{enumerate}
       \item train movement, continuous sample:
\begin{verbatim}
time object acceleration velocity distance track_from track_to dist_signal
\end{verbatim}
       \item discrete element event:
\begin{verbatim}
time object message
\end{verbatim}
      \end{enumerate}
       \insertplot{shunting}{Graph of train velocity during passage through shunting loop}
      \verb|distance| is total distance traveled by train, \verb|track_from| track where train beginning is located,
      \verb|track_to| track where train end is located, \verb|dist_signal| distance to nearest signal.

      In graph in figure \ref{shunting} is shown how train in model reacts to signal aspects
      during event of arrival at signal.
      Here one can see how control logic fails to allow path to following block in time,
      and train accelerates and stops between signals.
      In section \ref{trainMoveSection} is explained train behavior.

\chapter{Conclusion}
In this chapter are described possible program extensions and evaluation of entire thesis.
\section{Suggestions for Extensions}\label{chapExtensions}
There are many ideas how to extend current application, but here I will focus rather on those, for which
current implementation is directly designed, or which follow from it and are realizable in near future.

  \subsection*{Goals of Final Application}
    These goals I set out because I don't want implemented code to be created for only one purpose.
    By this I would like to somehow express possibilities of further development of entire application.
    Other students can build on my work and realize them.
    \begin{itemize}
        \item graphical editor and railway network control simulator in one application,
              focus on functionality, minimization of clicking
 	\item platform independence -- at least at source code level
 	\item data format that allows creating element templates by composition from others
 %	\item Automate routines in control process using fuzzy logic, neural networks or
 %   genetic algorithms. Not everything however will be possible immediately.
    \end{itemize}

  \subsection*{Physical Model of Movement}
    For current movement description also analytical solution exists. Description using differential equation
    however allows its change for more complex one, which will more faithfully model other possible cases.
    Train path (mainly in inter-station sections) I would conceive as curve
    in three-dimensional space -- i.e. another edge property.
    In doing so it is necessary to decompose current single force giving resulting acceleration into several components
    -- inclined plane on hills, friction, motor traction force, which could be controlled for example by fuzzy
    regulator. I can then also consider failures, during which motor force is zero.

  \subsection*{Automatic Block}\label{autoblokSection}
     Automatic block (AB) is permanently set, automatically renewed, path for multiple trains.
     It can be fixed \sourceRef{TrackBlock} assembled from several sections \sourceRef{SimpleTrack}.
     AB as whole has direction set by dispatcher. Individual sections
     immediately reserve path in that direction when their release occurs.
     In figure \ref{autoblok} is shown behavior of AB section.
     \insertimage{autoblok}{C/E Petri net of section in automatic block}
     AB should also be possible to create from set path.
     Furthermore with inter-station sections is related also use of other signal types (depends where they will be placed in network):
     \begin{description}
       \item[Section signal] in automatic block
       \item[Distant signal] using constant signal with \uv{Clear}
     \end{description}
  \subsection*{Saving Simulation State}
   Discrete process is fragmented into event commands, such as waiting, leading to activity loss. If by some universal way I create possibility to remember
   state in which each process is just passive between events and define transformation from beginning to this state, I can save this information together with other properties during serialization.
   Process could be internally represented as automaton or Petri net and its resp. its state would be
   saved and restored.
   For continuous variables it is necessary to save \sourceRef{jDisco.Variable}.
   It is possible that due to universal solution, it will require modifications in \emph{jDisco} library for serialization support. State of \sourceRef{LoopProcess} descendants, which inside iteration doesn't come to passive state,
   is possible to save already.
 \subsection*{Visualization Using Cellular Automata, Real-time}
   Square network of model directly predetermines using cellular automata for rendering control.
   The effort is to minimize amount of redrawn area and unnecessary computational load,
   caused by traversing states of all cells, by limitation only to small amount of cells of some types
   and for others I will proceed from assumption that if one cell changes, then probability
   of change in cells in neighborhood\footnote{non-empty cells} is set to maximum and then decreases with time.
   Furthermore it will be suitable to create \sourceRef{Monitor} with synchronization of real time with model time. It should
   better detect and suitably insert \uv{pauses} between continuous computation. Currently I handle this with discrete
   process, which for each second of model time, during which it waits i.e. other processes are simulated,
   calculates elapsed time and inserts directly needed pause of Java thread.
   Quasi-parallelism ensures that computation of entire simulation is delayed. It is thus insertion of some
   synchronization pulses, which however need to be scheduled at better moment.
 \subsection*{Interface to Model for Dispatcher Control}
   It is possible to build on current path formalization, implement their search.
   Create package of routines and functions covering access to model for control from several levels.
   This also follows completing commands from GUI interface.
   It will be suitable to formalize command sequence as language, which can be possibly optimized.

\section{Evaluation of Achieved Results}
Program simulates movement of several trains in network simultaneously.
Each train moves according to path determined during simulation -- reacts to signal aspects.
Another train cannot enter occupied tracks.
Established minimal goals were thus achieved, which example demonstrates. More complex network is reasonably demonstrable only in case of implementing control module and simulation visualization, when control must be taken over by human.
Validity of this simulation model could be verified by collecting statistics about real train networks and comparing
with model.

Suitable features of languages resp. tools finally prevail, although some constructs in them were more difficult
to express.
If I had to verify what load simulation is on operating system, always values during
running virtual machine instance\footnote{I ran shunting loop without synchronization with real time for 15 minutes} moved on my device around these values and didn't change significantly:
shared memory -- 10\,MB, code and data in physical memory -- 25\,MB, amount of virtual memory -- 213\,MB,
CPU performance 91\,\% (this is not attempt at sophisticated performance analysis, I only monitored whether program doesn't demand unnecessarily many resources).

On own initiative I tried possibility to work on such project and
immerse myself in problematics of analysis and modeling of complex systems,
which includes dealing with case studies of technical devices.
Furthermore I expanded horizon about how to incorporate simulation,
and improved in application design and programming.
Problem I had especially at beginning with estimation of time complexity of work.
% numerical error
Coding itself went quite quickly, but I forgot about debugging in time estimates.
I prefer to test continuously, error is then easier to find. For extensive projects however it is still necessary to create
testing scripts.
In design area there is still what to learn and more utilize design patterns and mainly adhere to their principles.
But I must note that time spent studying and designing system was longer than programming time, which I would like in continuation of building this application to further shorten by avoiding mistakes which I committed.

\bibliographystyle{czechiso}
\begin{flushleft}
 \bibliography{bakalarka}
\end{flushleft}

\zkratky

\end{document}
